\section{Related work}
\label{sec:related}

The result sits at the intersection of three lines: the RLVR reasoning-boundary
debate, the theory of entropy and concentration in policy optimization, and the
sign-variation geometry of totally positive kernels. We review each, then mark
the boundary with the two concurrent manuscripts precisely.

\paragraph{RLVR algorithms and systems.}
Reinforcement learning with verifiable rewards grew out of policy-gradient and
RLHF pipelines \citep{schulman2017ppo,ouyang2022instructgpt,ahmadian2024rloo,hu2025reinforceplusplus,kazemnejad2024vineppo}
and was crystallized by GRPO and the DeepSeek line
\citep{shao2024deepseekmath,deepseekai2025r1,yu2025dapo,liu2025drgrpo}, the term
itself fixed in \citep{lambert2024tulu3}. Open ``zero'' recipes train from the
base model without supervised warmup
\citep{hu2025orz,zeng2025simplerl,zhao2025absolutezero,yan2025luffy,luong2024reft},
and large system reports scale the setup
\citep{kimiteam2025k15,he2025skywork,mistral2025magistral,bercovich2025llamanemotron,yang2024qwen25math,wei2025swerl,khatri2025scalerl}.
Rewards come from outcome checkers and process models
\citep{cobbe2021gsm8k,uesato2022process,lightman2023verify,wang2024mathshepherd,zhang2025prmlessons,ma2025generalreasoner};
recent surveys map the area \citep{xu2025surveyLRM,zhang2025surveyRL}. None states
a pass@$k$ crossing count or a per-prompt no-free-lunch condition.

\paragraph{The phenomenon and the debate.}
\citet{yue2025rlincentivize} report that RLVR raises pass@1 while the base model
overtakes it at large $k$ and read the reversal as reweighting rather than new
capability; \citet{liu2025prorl} argue prolonged RL expands the boundary and
keeps the trained model ahead; \citet{zhou2026shrink} diagnoses the inversion and
proposes a diversity intervention. A cluster of work reads RL as amplifying or
sharpening base behavior rather than adding it
\citep{zhao2025echochamber,setlur2024incorrect,huang2024sharpening,shao2025spurious,wu2025invisibleleash}.
Reward hacking and chain-of-thought obfuscation are documented failure modes of the same optimization \citep{baker2025monitoring}. Our Theorem \ref{thm:price} adjudicates: the crossover is forced when the loss
reaches a finely covered hard tail (Assumption \ref{ass:fine}) and absent
otherwise, so the Yue and ProRL observations are two sides of one conditional
statement.

\paragraph{pass@$k$, sampling, and test-time compute.}
The pass@$k$ metric originates with code evaluation \citep{chen2021codex}, and
sampling-based inference spans self-consistency and its universal form
\citep{wang2022selfconsistency,chen2023universalsc}, repeated sampling and
inference scaling \citep{brown2024monkeys,snell2024testtime,wu2024inferencescaling,muennighoff2025s1},
and best-of-$N$ alignment \citep{amini2024variationalbon,dong2023raft}. A recent
line optimizes pass@$k$ directly or debates it as an objective
\citep{walder2025passkppo,chen2025passktraining,yu2025passkdiagnostic,gui2025inferenceobj}.
Chain-of-thought and bootstrapped reasoning are the substrate
\citep{wei2022cot,zelikman2022star,yuan2023rft}. We do not optimize the metric;
we bound how the base-to-RL shift can move the curve and price the move.

\paragraph{Entropy, exploration, and diversity.}
Entropy collapse and its exploitation are central to RLVR dynamics
\citep{cui2025entropy,agarwal2025entropymin,wang2025entropytokens}, extending the
maximum-entropy tradition \citep{haarnoja2018sac}. Alignment is known to narrow
output diversity and induce mode collapse
\citep{kirk24diversity,mohammadi2024creativity,lanchantin2025divpo,zhang2025verbalized,jiang2025hivemind}.
These describe or counter concentration; the per-prompt accounting and the
crossing bound are new here.

\paragraph{RLHF and preference-optimization theory.}
The reward-model and preference-learning foundations
\citep{bradleyterry52,christiano17,stiennon20summarize,rafailov23dpo,azar24ipo,ethayarajh24kto,zhu23principled,munos24nash}
and the Goodhart limit of reward optimization \citep{gao23overopt} frame RL as a
one-step contextual bandit, the view we take. Our no-free-lunch duality is a
Pareto statement on the strategy simplex, which these do not address.

\paragraph{Policy-gradient theory.}
The tilt identity rests on policy-gradient and mirror-descent analysis
\citep{sutton99pg,kakade01npg,agarwal21pgtheory,tomar22mdpo,mei20softmax,schulman15trpo},
control-as-inference \citep{levine18control}, and the bandit view
\citep{lattimore20bandit}.

\paragraph{Named-phenomenon theory, for positioning.}
Pure-theory explanations of named LLM phenomena are the archetype this paper
follows: hallucination
\citep{kalai25hallucinate,xu24hallinevitable}, superposition and in-context
learning \citep{elhage22superposition,vonoswald23icl,olsson22induction}, grokking
\citep{power22grokking,nanda23progress}, emergence and its critique
\citep{wei22emergent,schaeffer23mirage}, and positional and attention effects
\citep{liu24lostmiddle,xiao24attnsink}.

\paragraph{Mathematical tools.}
The crossing bound is a variation-diminishing property of totally positive
kernels \citep{karlin68tp,schoenberg51pff,schoenberg30vdt,karlinstudden66}, with
the exponential-sum sign count going back to Descartes' rule
\citep{jameson06descartes} and Laplace-transform sign changes
\citep{widder41laplace}; the no-free-lunch condition is a Pareto statement proved
by a theorem of the alternative \citep{mangasarian69nlp,miettinen99moo}, and the
frontier uses stochastic dominance \citep{shaked07stochorders}. The novelty is
the target, not the tools: pricing pass@1 against per-prompt success and bounding
pass@$k$ crossings by $\Svar(G)$.

\section{Concurrent work}
\label{sec:concurrent}

Two manuscripts from the last two months are the nearest neighbors; the overlap
is partial and real.

\paragraph{Yang and Chen, ``RLVR is a Kernel, Not a Function''
\citep{yangchen2026kernel}.}
They build statistical inference for pass@$k$ crossovers and prove, in the
\emph{function} case with zero latent-failure mass, a single interior fixed
point, and a side-preserving map, that the advantage changes sign at most once
(their Theorem 3). They note multiple crossings become possible once the map is a
kernel or the failure mass is positive, and leave that case without a bound
(their Remark 1), framing two supported sign changes as a falsification of the
single-fixed-point class. Our Theorem \ref{thm:price}(a) supplies the missing
bound: for \emph{any} shift the number of crossings is at most $\Svar(G)$,
kernels included, and it is tight; their single-crossing theorem is our
order-preserving corollary \ref{thm:price}(d). Our no-free-lunch result (Theorem
\ref{thm:nfl}) and the price frontier (Theorem \ref{thm:frontier}) have no
counterpart in their work. Their prompts are independent given a common map; our
strategies are shared across prompts, the microfoundation of the kernel
dispersion they posit.

\paragraph{``Winner-Take-All'' (ICLR 2027 submission S4R2zYtFhs).}
They study the entropy side of the same tilt in a finite-mode, single-prompt
mean-field model: the exact entropy-mass curve, the shifted-exponential as a
finite-window approximation with error bounds, a proof that the terminal law is
not exponential, and the collapse of correct-mode entropy to a single winner. We
claim none of these. Our entropy statements (Theorem \ref{thm:tiltlaw}, the
lemmas) are the strategy-level, mesoscopic context that links their entropy side
to our pass@$k$ side through one tilt identity; where they overlap, priority is
theirs. Their model has one prompt and many modes; ours has many prompts sharing
strategies, which is what makes pass@$k$ well defined.

\paragraph{Summary.}
The crossover's statistics \citep{yangchen2026kernel} and the entropy trajectory
(Winner-Take-All) are concurrent. This paper adds the general crossing bound with
tightness, the per-prompt no-free-lunch condition with its Farkas exchange rate,
the forced-crossover result under fine coverage with its counterexample, and the
price frontier. The unifying object is the shared-strategy tilt.